% Introduction and Related Work draft (Stage 2 output).
% All \cite{} keys are defined in references.bib.

\section{Introduction}

Higher-order thalamic nuclei were once viewed as passive relays of cortical
activity, but a growing body of work has repositioned them as active
participants in cognition and motor control \cite{sherman2007thalamic,
sherman2016thalamus, halassa2019thalamocortical}. Unlike primary sensory
thalamic nuclei, higher-order nuclei receive the bulk of their drive from
cortex itself, forming reciprocal cortico-thalamo-cortical (CTC) loops that
have been mapped in rodent sensory, motor, and prefrontal regions
\cite{yamawaki2015synaptic, economo2018distinct, collins2018reciprocal,
guo2018anterolateral}. Perturbation studies in mice demonstrate that these
loops are required for the maintenance of persistent activity during working
memory \cite{guo2017maintenance, bolkan2017thalamic, schmitt2017thalamic},
for the execution of dexterous movements \cite{sauerbrei2020cortical}, and
for flexible switching among cortical representations
\cite{rikhye2018thalamic, alcaraz2018thalamocortical}.

The central computational puzzle raised by this anatomy is that thalamic
populations are much smaller than the cortical populations with which they
interact. A CTC loop therefore implements a compression of cortical
activity into a low-dimensional thalamic space, followed by an expansion
back to cortex. For purely feedforward architectures, theories of
compressed sensing and efficient coding have clarified how such bottlenecks
should be structured to preserve or emphasize behaviorally relevant signals
\cite{ganguli2012compressed, olshausen1996emergence}. Recurrent CTC loops,
however, are not feedforward; insights obtained in feedforward settings do
not transfer directly. Existing theoretical treatments of thalamocortical
loops tune the weights of the loop by numerical optimization and show that
the resulting circuit can support motor preparation or flexible motor
sequencing \cite{kao2021optimal, logiaco2021thalamic}. What these models do
not address is how such tuned weights could be established by a
biologically plausible mechanism.

The question becomes especially pointed when we ask how local, error-driven
plasticity could operate in a CTC loop. Three-factor and random-feedback
rules have been proposed as biologically plausible approximations to
gradient-based learning in recurrent networks \cite{murray2019local,
lillicrap2016random, gilra2017predicting, miconi2017biologically,
gerstner2018eligibility}, and experimental studies support plasticity at
thalamocortical synapses into adulthood
\cite{hasegawa2018thalamocortical}. Yet evidence for plasticity at
corticothalamic synapses is limited, and these synapses sit far upstream of
the behavioral readout so that any local learning rule must contend with a
particularly severe credit-assignment problem. We hypothesized that the
ability of the thalamus to contribute to learning therefore depends on the
\emph{structure} of corticothalamic projections, rather than on local
plasticity at those projections, and asked what structure is optimal for
each task.

We approach this question with a recurrent thalamocortical model in which
$N$ cortical neurons are reciprocally connected to $M$ uncoupled thalamic
neurons, with $M \ll N$. Thalamocortical weights are updated with a local
three-factor rule (RFLO) \cite{murray2019local, lillicrap2016random},
while corticothalamic weights are optimized either by meta-learning on a
slow outer-loop timescale \cite{wang2018prefrontal, hospedales2022meta} or
by idealized alignment with specific cortical subspaces. Our central
finding is that the optimal corticothalamic structure is strongly
task-dependent: autonomous motor control benefits most from a projection
carrying an efference copy of the readout, while working memory benefits
most from a projection carrying the leading principal components of
cortical activity. Reanalysis of intracortical and thalamic recordings
from mice performing grasping and delayed-discrimination tasks
\cite{sauerbrei2020cortical, guo2017maintenance} is consistent with this
prediction. Our contributions are:

\begin{itemize}
  \item We show that a local three-factor plasticity rule succeeds at
  thalamocortical synapses but fails at corticothalamic synapses, because
  feedback alignment does not propagate through the thalamic bottleneck.
  \item We demonstrate that meta-learning identifies task-specific
  corticothalamic structures that reduce task error by a median factor of
  5.8 (motor control) and 5.1 (working memory) relative to random
  projections.
  \item We show that idealized readout-aligned projections are optimal for
  motor control while principal-component-aligned projections are optimal
  for working memory, and that these choices generalize to a composite
  delayed-reaching task in a model with separate preparatory and
  execution thalamic modules.
  \item We reanalyze published neural recordings from motor cortex /
  motor thalamus (3 sessions, 101 thalamic neurons) and from ALM /
  VM--VAL thalamus (5 sessions, 72 thalamic neurons) and find that the
  cortical subspaces captured by thalamus are consistent with the
  task-specific structures predicted by the model.
\end{itemize}

\section{Related Work}

\paragraph{Higher-order thalamus and cortico-thalamo-cortical loops.}
Reciprocal CTC loops are a frequent anatomical motif across rodent
cortical areas \cite{yamawaki2015synaptic, economo2018distinct,
collins2018reciprocal, guo2018anterolateral, halassa2019thalamocortical}.
Inactivation and optogenetic experiments show that these loops are
required for the maintenance of persistent activity during working
memory and for the execution of prepared movements
\cite{guo2017maintenance, bolkan2017thalamic, schmitt2017thalamic,
sauerbrei2020cortical, rikhye2018thalamic, alcaraz2018thalamocortical}.
Theoretical models have shown that tuned CTC weights can support motor
preparation and flexible motor sequencing
\cite{kao2021optimal, logiaco2021thalamic}, but they rely on numerical
optimization and do not specify a biological mechanism for acquiring the
weights. Our work fills this gap by supplying structured corticothalamic
connectivity and a local plasticity rule that together realize
biologically plausible learning.

\paragraph{Biologically plausible learning in recurrent networks.}
Classical recurrent learning is non-local
\cite{williams1989learning}. To recover locality, a number of
approximations have been introduced: random feedback alignment
\cite{lillicrap2016random}, three-factor eligibility rules
\cite{gerstner2018eligibility, bittner2017behavioral}, dendritic
segregation for deep credit assignment \cite{guerguiev2017towards},
stable local learning in spiking RNNs \cite{gilra2017predicting},
biologically plausible RNN training reproducing task dynamics
\cite{miconi2017biologically}, and the specific Random-Feedback-Local-Online
(RFLO) rule that we use here \cite{murray2019local}. These rules have
been analyzed for synapses that directly drive the readout; our results
address the less-studied regime of synapses that are two steps removed
from the readout through the thalamic bottleneck.

\paragraph{Low-rank recurrent network theory and meta-learning.}
Recent theoretical work conceptualizes recurrent connectivity as the sum
of a random full-rank component and a tunable low-rank perturbation
\cite{mastrogiuseppe2018linking}. Because the CTC loop contributes a
rank-$M$ perturbation to the effective cortical connectivity, our
analysis inherits these tools. We combine this framework with
meta-learning \cite{wang2018prefrontal, hospedales2022meta}: a slow outer
loop optimizes the corticothalamic weights, while a fast inner loop uses
a local plasticity rule to update thalamocortical weights. This two-timescale
decomposition is consistent with the view that corticothalamic
structure reflects evolution or development, while thalamocortical
synapses are plastic throughout life.

\paragraph{Cortical dynamics during motor learning and cognition.}
Recurrent dynamics have been shown to support flexible computation in
prefrontal cortex and to reorganize during motor learning in rodent motor
cortex \cite{mante2013context, peters2014emergence,
churchland2012neural}. These studies largely attribute learning-induced
changes to plasticity within cortex; our results argue that at least some
of the learning signal is routed through thalamus, and that the ability
of the thalamus to participate depends on the specific structure of
corticothalamic projections.
